This work presented a new perspective on automatic figure captioning, demonstrating that a language-based approach, \ie, summarizing figure-referring paragraphs, can outperform conventional vision-based methods.
Our analysis further showed many unhelpful captions in arXiv papers, highlighting data quality's impact on captioning performance.
%
This work lays the groundwork for further research, including exploring new data selection, revision, and augmentation strategies to mitigate the effects of low-quality data, developing new evaluation methods, and creating more robust models that better handle noisy data.
%Additionally, we plan to extend the application of the proposed technology to cover a broader range of figures and article types.
We also aim to expand the technology's scope to cover a wider variety of figures and article types.

%--------------- dead kitten -----------

\begin{comment}



This paper introduced a novel approach and perspective to automatic figure captioning. 
We showed that a language approach, \ie, summarizing the paragraphs that explicitly refer to a figure, can more effectively produce high-quality captions for that figure than traditional vision approaches. 
Our analysis also showed that a significant number of captions in arXiv \texttt{cs.CL} papers were considered unhelpful to readers.
These findings collectively motivate unexplored new possibilities. 
Examples include new data selection, revision, or augmentation strategies that mitigate the effects of low-quality data on model performance; new evaluation approaches beyond using the author-written caption as gold; and new, robust models that work better with noisy data. 
We will develop automatic post-editing models that spot and fix incorrect information in generated captions. We will also adopt the proposed technology for more types of figures and articles.

This paper introduced a novel approach to automatic figure captioning. 
We showed that summarizing the paragraphs that explicitly refer to a figure 
can effectively produce high-quality captions for that figure. 
Our analysis also showed that a significant number of captions in arXiv \texttt{cs.CL} papers were considered unhelpful to readers. 

As our paper pointed out that 50\% of the captions are bad, 
it motivates a set of unexplored plausible future directions, including 
(\textit{i}) automatic data selection/filtering, 
(\textit{ii}) data revision, 
(\textit{iii}) new evaluation methods beyond treating the one author-written caption as gold, and 
(\textit{iv}) new modeling techniques that work with noisy data. 
In the future, we will improve captioning performance 
by automatically spotting and fixing incorrect information in generated captions. 
We will also adopt the proposed technology to a broader type of figures and articles.

\kenneth{This can be longer. Say more about what we're going to do next.}

\end{comment}